\documentclass[tcc,capa]{texufpel}

\usepackage[utf8]{inputenc} % acentuacao
\usepackage{graphicx} % para inserir figuras
\usepackage[T1]{fontenc}

\hypersetup{
    hidelinks, % Remove coloração e caixas
    unicode=true,   %Permite acentuação no bookmark
    linktoc=all %Habilita link no nome e página do sumário
}

\unidade{Centro de Desenvolvimento Tecnológico}
\curso{Ciência da Computação}
\nomecurso{Bacharelado em Ciência da Computação}
\titulocurso{Bacharel em Ciência da Computação}

\unidadeeng{Technology Development Center}
\cursoeng{Computer Science}


\title{Gerenciamento Multi-Estruturado de Dados Móveis Gamificados: Integrando Fontes Heterogêneas.}

\author{Santos}{Rávilon Aguiar dos}
\advisor[Prof.~Dr.]{Fleischmann}{Ana Marilza Pernas}

%Palavras-chave em PT_BR
%use letras minúsculas seguidas por ;
%a última terminada por .
\keyword{Banco de Dados Não Relacional;}
\keyword{Integração com dispositivos Móveis;}
\keyword{Fontes Heterogêneas;}
\keyword{Gerenciamento de Banco de Dados.}

%Palavras-chave em EN_US
\keywordeng{Non-Relational Database;}
\keywordeng{Integration with Mobile Devices;}
\keywordeng{Heterogeneous Sources;}
\keywordeng{Database Management.}

\begin{document}

%\renewcommand{\advisorname}{Orientadora}           %descomente caso tenhas orientadora
%\renewcommand{\coadvisorname}{Coorientadora}      %descomente caso tenhas coorientadora

\maketitle 

\sloppy

\fichacatalografica

%\folhadeaprovacao

%Composição da Banca Examinadora

%Resumo em Portugues (no maximo 500 palavras)
\begin{abstract}
Este Trabalho de Conclusão de Curso, intitulado "Gerenciamento Multi-Estruturado de Dados Móveis Gamificados: Integrando Fontes Heterogêneas", propõe o desenvolvimento de um modelo e arquitetura de bancos de dados para a gestão eficiente de dados provenientes de fontes heterogêneas. O objetivo principal é suportar uma aplicação que incentiva a redução do consumo energético através da gamificação, focando no usuário desse ambiente.

A proposta utiliza bancos de dados NoSQL para capturar e armazenar dados de forma periódica em uma edificação pública, processando essas informações com princípios de gamificação. O trabalho aborda a modelagem da estrutura dos dados provenientes de diferentes fontes, como usuários, sensores em salas e corredores, e tecnologias disponíveis. A estrutura multi-modelo é apresentada como a mais adequada para cenários de gamificação devido à necessidade de integrar e adaptar dados de diversas origens.

O estudo considera cenários e requisitos específicos de uma aplicação de gamificação, com base em casos de estudo do projeto "Validação de sistema de gamificação, gestão e automação na nuvem para eficiência energética de edificações públicas com foco inicial em universidades". Este projeto, realizado na Universidade Federal de Pelotas (UFPel), visa identificar o impacto da gamificação na economia de energia e no Desempenho Energético Operacional (DEO) de uma edificação, utilizando como objeto de estudo o prédio do Centro de Artes da UFPel.
\end{abstract}

%Resumo em Inglês (no maximo 500 palavras)
\begin{englishabstract}{Multi-Structured Gamified Mobile Data Management: Integrating Heterogeneous Sources}
This Thesis, titled "Multi-Structured Management of Gamified Mobile Data: Integrating Heterogeneous Sources," proposes the development of a model and database architecture for the efficient management of data from heterogeneous sources. The main goal is to support an application that encourages energy consumption reduction through gamification, focusing on the user of this environment.

The proposal uses NoSQL databases to periodically capture and store data in a public building, processing this information with gamification principles. The work addresses the modeling of data structures from different sources, such as users, sensors in rooms and corridors, and available technologies. The multi-model structure is presented as the most suitable for gamification scenarios due to the need to integrate and adapt data from various sources.

The study considers specific scenarios and requirements for a gamification application, based on case studies from the project "Validation of a gamification system, management, and cloud automation for energy efficiency in public buildings with an initial focus on universities." This project, carried out at the Federal University of Pelotas (UFPel), aims to identify the impact of gamification on energy savings and Operational Energy Performance (OEP) of a building, using the UFPel Arts Center as the study object.
\end{englishabstract}

%Lista de Figuras
\listoffigures

%Lista de Tabelas
\listoftables

%lista de abreviaturas e siglas
\begin{listofabbrv}{NoSQL}%coloque aqui a maior sigla para ajustar a distância
        \item[\textbf{IoT}] Internet of Things (Internet das coisas)
        \item [\textbf{NoSQL}] Não relacional
  
\end{listofabbrv}

%Sumario
\tableofcontents

\chapter{Introdução}

\section{Motivação}

A gamificação tem se mostrado uma ferramenta poderosa para aumentar o engajamento e motivação dos usuários em diversos contextos, desde a educação até a gestão de saúde. Gamificação, conforme descrito por Werbach e Hunter (2012), é "a aplicação de elementos de jogos e técnicas de design de jogos em contextos que não são jogos" ~\cite{werbach2012}. Segundo os mesmos autores, objetiva-se com a gamificação aplicar mecânicas de jogo e design de experiências para envolver digitalmente os usuários e motivá-los na obtenção de seus objetivos.

No caso do projeto “Validação de sistema de gamificação, gestão e automação na nuvem para eficiência energética de edificações públicas com foco inicial em universidades”, alvo deste trabalho, busca-se demonstrar que técnicas de gamificação têm o potencial de incentivar comportamentos sustentáveis e reduzir o consumo de energia por meio de um sistema de recompensas e feedback contínuo. A proposta dos autores é criar um sistema que estimule os servidores públicos a adotarem posturas de sustentabilidade, contribuir para uma gestão mais efetiva e consequente redução do consumo energético de edificações públicas. Para isso, serão implementados sistemas de medição e verificação do consumo energético no prédio do Centro de Artes da UFPel, envolvendo software e hardware de gestão de energia.

Visando contribuir para este projeto, o presente Trabalho de Conclusão de Curso foca nos aspectos relacionados a modelagem dos dados sensorados nestes ambientes, assim como seu correto armazenamento e manipulação viabilizando posteriores desenvolvimentos de aplicações a partir deste.

Neste aspecto, a integração de sensores em ambientes diversos, como edifícios, cidades e dispositivos móveis, é foco de pesquisa devido ao desafio de se gerenciar, de maneira eficiente, uma vasta quantidade de dados heterogêneos.

Tradicionalmente, os bancos de dados relacionais eram a principal solução para armazenamento e processamento de dados. No entanto, a rigidez de seu modelo de dados e a dificuldade de escalar horizontalmente tornaram-se limitações significativas no contexto de aplicações mais modernas, como Big Data ou Internet das Coisas, onde os dados são altamente variados e crescem rapidamente. Nesse cenário, destacam-se os bancos de dados NoSQL, que oferecem maior flexibilidade e escalabilidade para lidar com grandes volumes de dados não estruturados e semi-estruturados \cite{Stonebraker2005, Grolinger2013}.

De acordo com Sadalage e Fowler (2012), NoSQL é "uma abordagem moderna para o gerenciamento de dados que adota uma estrutura de dados mais flexível e escalável do que a oferecida pelos bancos de dados relacionais tradicionais" \cite{Sadalage2012}. Os bancos de dados NoSQL permitem armazenar e processar dados de maneira eficiente, suportando a integração de fontes heterogêneas de informação. De acordo com Sadalage e Fowler (2012, p. 12), "a escolha de uma solução NoSQL deve ser baseada em uma compreensão clara das necessidades específicas da aplicação" \cite{Sadalage2012}. Esta flexibilidade é crucial para o gerenciamento de dados de sensores, onde os formatos de dados podem variar amplamente.

A combinação dessas duas abordagens – a gestão de dados por bancos de dados NoSQL e sua disponibilidade para fins de gamificação – fundamenta a proposta deste Trabalho de Conclusão de Curso. Este projeto visa desenvolver um modelo e uma arquitetura de banco de dados que suporte a coleta e o processamento eficiente de dados, integrando princípios de gamificação para incentivar a redução do consumo energético em um edifício específico. Em relação aos dados, é necessário tratar três tipos principais de informações \cite{Sagl2015}:

\begin{enumerate}
    \item Dados de Sensores Ambientais (Environmental Sensing): Esses dados são coletados por sensores instalados no edifício que monitoram o ambiente interno, fornecendo informações cruciais sobre fatores como temperatura, umidade e qualidade do ar, e consumo energético do ambiente.
    \item Dados de Sensores Móveis (Mobile Sensors): Estes são dados provenientes da aplicação de gamificação, onde dispositivos móveis dos ocupantes do edifício registram atividades e comportamentos relacionados ao consumo energético.
    \item Dados de Pessoas como Sensores (People as Sensors): Referem-se aos dados fornecidos diretamente pelos ocupantes do edifício, muitas vezes através de interações com a aplicação de gamificação, como feedback sobre o ambiente ou participação em desafios de economia de energia.
\end{enumerate}

Esses dados contextualizados são essenciais para criar um sistema inteligente capaz de interferir no funcionamento da cidade. A integração eficaz desses diversos tipos de dados permitirá não só a monitorização precisa do consumo energético, mas também a implementação de estratégias de gamificação que motivem comportamentos sustentáveis entre os ocupantes do edifício. Definir esses dados, no contexto de aplicação do projeto, consiste em uma das etapas do trabalho a ser desenvolvido.

Espera-se disponibilizar à aplicação de gamificação, para acesso e manipulação, os dados presentes na arquitetura por meio de uma Interface de Programa de Aplicação (do inglês, Application Program Interface - API). De acordo com Goodwin (2024), uma API consiste de um conjunto de regras ou protocolos que permite que aplicações de software comuniquem entre si para trocar dados, funcionalidades e recursos \cite{Goodwin2024}. APIs simplificam e aceleram o desenvolvimento de software, permitindo a integração de dados e serviços de outras aplicações sem a necessidade de desenvolvimento do zero, ademais não restringindo esses dados apenas a seu objetivo inicial, voltado ao suporte de uma aplicação para a redução do consumo energético.

Essa proposta também se alinha com iniciativas globais de sustentabilidade e eficiência energética. A necessidade de reduzir o consumo de energia e promover práticas sustentáveis está em destaque em agendas de políticas públicas e privadas. Tal ênfase na urgência da redução de consumo energético é possível de ser notada por meio da Agenda 2030, das Nações Unidas no Brasil (2015) \cite{agenda2030} a qual cita a urgência de se investir em tecnologia para possibilitar a redução no consumo energético. O caso de uso específico, focado no prédio do Centro de Artes da Universidade Federal de Pelotas, é um exemplo concreto de como a tecnologia pode ser utilizada para atingir esses objetivos. A implementação de sensores para monitorar o consumo energético e a utilização de uma aplicação gamificada para incentivar a economia de energia pode servir de modelo para outras instituições e edificações públicas.

Distinguindo-se de outros trabalhos já realizados, este projeto se destaca pela sua abordagem integrada que combina a flexibilidade e escalabilidade dos bancos de dados NoSQL com a capacidade de engajamento da gamificação. Conforme Sadalage e Fowler (2012), "os sistemas NoSQL são particularmente eficazes para aplicações que necessitam de alta disponibilidade e escalabilidade horizontal" \cite{Sadalage2012}. Enquanto muitos projetos focam apenas em um desses aspectos, a integração proposta oferece uma solução holística que aborda tanto a gestão eficiente de dados quanto a motivação dos usuários para comportamentos sustentáveis. Isso representa um grau de inovação significativo, com potencial para influenciar futuras pesquisas e práticas em gestão de dados e sustentabilidade.

Em conclusão, a proposta "Gerenciamento Multi-Estruturado de Dados Móveis Gamificados: Integrando Fontes Heterogêneas" não apenas responde às necessidades atuais de gestão eficiente de dados e redução de consumo energético, mas também propõe uma abordagem inovadora que pode ter impacto duradouro na academia e na comunidade.

\chapter{Principais Conceitos}

\section{O que é Gamificação}

Gamificação refere-se ao uso de elementos e técnicas de design de jogos em contextos não lúdicos com o objetivo de engajar usuários e resolver problemas do mundo real. Segundo Deterding et al. e discutido por Kim \cite{kim2015gamification}, a gamificação se distingue dos jogos completos pelo seu foco em aplicar mecânicas de jogo fora de um ambiente puramente lúdico. Zichermann e Cunningham definem gamificação como "o processo de usar o pensamento de jogos e mecânicas de jogos para engajar usuários e resolver problemas."

Aplicações práticas de gamificação incluem o Nike+, que usa elementos de jogo para motivar a prática de exercícios físicos, e o Chore Wars, que transforma tarefas domésticas em desafios para aumentar a produtividade. Esses exemplos demonstram a eficácia da gamificação na mudança de comportamento e na resolução de problemas específicos.

Deterding et al. também discutem a distinção entre gamificação e design lúdico. Enquanto a gamificação utiliza elementos de jogos estruturados para resolver problemas reais, o design lúdico emprega um uso mais expressivo e improvisacional de comportamentos, sem regras rígidas ou competição direta. Em contraste com jogos tradicionais, onde o objetivo é alcançar metas fictícias, a gamificação busca resolver problemas reais e promover comportamentos desejados.

Em resumo, a gamificação é uma abordagem que utiliza características dos jogos para incentivar a participação ativa, motivar comportamentos desejados e resolver problemas reais, diferenciando-se tanto dos jogos quanto do design lúdico na aplicação de seus elementos \cite{kim2015gamification}.

\section{Metodologias de Gamificação}

\subsection{Elementos e Mecânicas de Jogos}

Metodologias de gamificação frequentemente aplicam elementos de jogos, como pontos, medalhas, rankings e desafios, em contextos não relacionados a jogos. Werbach e Hunter destacam a importância do design estratégico e da incorporação desses elementos para atingir objetivos específicos, como aumentar a participação ou melhorar os resultados de aprendizagem \cite{werbach2012for}.

\subsection{Design Motivacional}

A compreensão das motivações intrínsecas e extrínsecas é fundamental na gamificação. Estudos ressaltam a importância de adaptar as mecânicas de jogos para aprimorar a motivação intrínseca, promovendo um senso de autonomia, domínio e propósito, em vez de depender exclusivamente de recompensas externas, como pontos \cite{ryan2000intrinsic}.

\subsection{Contextos Educacionais}

A gamificação tem sido amplamente estudada em ambientes educacionais, envolvendo o uso de técnicas de aprendizagem baseadas em jogos (GBL) para melhorar o engajamento, desempenho e motivação dos estudantes. Além disso, diferentes metodologias exploram como a gamificação pode desenvolver habilidades interpessoais e uma mentalidade empreendedora entre os alunos \cite{kapp2012gamification}.

\subsection{Adaptação ao Contexto}

Metodologias eficazes de gamificação consideram o contexto em que são aplicadas. Por exemplo, durante a pandemia de COVID-19, educadores utilizaram a gamificação para ajudar os alunos a se adaptarem ao ensino online. A integração da gamificação na educação de liderança e no treinamento vocacional demonstra como a adaptação contextual dos elementos de jogos pode gerar melhores resultados \cite{deterding2011game}.

\subsection{Desafios e Limitações}

As metodologias de gamificação devem também abordar desafios como a necessidade de recursos adequados, o risco de experiências mal projetadas e a dificuldade de manter o engajamento ao longo do tempo \cite{deterding2011game}.

\section{Definição dos Dados do Usuário - Contexto da Aplicação}

Neste capítulo, detalhamos os diferentes tipos de dados gerados pelos usuários e sua relação com a aplicação de gamificação voltada para a eficiência energética em edificações públicas. Compreender esses dados é essencial para desenvolver um sistema que não apenas monitore e analise o consumo de energia, mas também engaje os usuários em comportamentos sustentáveis. A seguir, são discutidos os principais tipos de dados a serem considerados.

\subsection{Dados de Sensores Ambientais}

\textbf{Descrição:} Dados coletados através de sensores instalados no ambiente, como no prédio do Centro de Artes da UFPel. Esses sensores monitoram variáveis ambientais essenciais para a gestão eficiente da energia.

\textbf{Exemplos de Dados:}
\begin{itemize}
    \item Temperatura interna do ambiente
    \item Umidade relativa
    \item Qualidade do ar
    \item Consumo energético do ambiente
\end{itemize}

\textbf{Relevância para a Aplicação:} Esses dados fornecem uma visão detalhada das condições ambientais e do consumo energético, permitindo identificar padrões e anomalias. Informações precisas são fundamentais para o feedback fornecido aos usuários e para a implementação de estratégias de economia de energia.

\subsection{Dados de Sensores Móveis}

\textbf{Descrição:} Dados provenientes dos dispositivos móveis dos usuários coletados através da aplicação de gamificação. Incluem informações sobre atividades e comportamentos dos usuários em relação ao consumo de energia.

\textbf{Exemplos de Dados:}
\begin{itemize}
    \item Atividades registradas pelos dispositivos móveis
    \item Comportamentos relacionados ao consumo energético
    \item Interações com desafios e recompensas na aplicação
\end{itemize}

\textbf{Relevância para a Aplicação:} Esses dados ajudam a entender como os comportamentos dos usuários impactam o consumo de energia e a ajustar as mecânicas de gamificação para promover práticas mais sustentáveis. A análise desses dados também personaliza o feedback e as recompensas oferecidas.

\subsection{Dados de Pessoas como Sensores}

\textbf{Descrição:} Dados fornecidos diretamente pelos usuários através da aplicação de gamificação. Incluem feedback sobre o ambiente e participação em atividades de economia de energia.

\textbf{Exemplos de Dados:}
\begin{itemize}
    \item Feedback sobre condições ambientais e sugestões de melhorias
    \item Participação em desafios e campanhas de economia de energia
    \item Respostas a pesquisas e questionários sobre práticas de consumo energético
\end{itemize}

\textbf{Relevância para a Aplicação:} Esses dados capturam a percepção dos usuários sobre o ambiente e as iniciativas de gamificação. A interação direta com os usuários permite ajustar a aplicação e as estratégias de gamificação, aumentando a eficácia das ações voltadas para a redução do consumo energético.

\subsection{Integração e Análise dos Dados}

\textbf{Descrição:} A integração dos diferentes tipos de dados é crucial para criar um sistema inteligente que analise e utilize essas informações de forma coesa. A análise dos dados permitirá identificar padrões, medir o impacto das estratégias de gamificação e fazer ajustes conforme necessário.

\textbf{Abordagem:}
\begin{itemize}
    \item Coleta e armazenamento dos dados em um banco de dados NoSQL para garantir flexibilidade e escalabilidade
    \item Processamento e análise dos dados para gerar insights sobre o consumo energético e o engajamento dos usuários
    \item Desenvolvimento de relatórios e dashboards para visualização dos dados e acompanhamento das métricas de desempenho
\end{itemize}

\textbf{Relevância para a Aplicação:} A integração e análise dos dados permitirão uma visão abrangente do impacto das ações de gamificação e da eficiência energética no prédio. Isso facilitará a implementação de melhorias contínuas e a adaptação das estratégias para maximizar os resultados.

\section{Bancos de Dados Voltados para Gamificação}

\subsection{Natureza dos Dados a Serem Armazenados}

Em aplicações de gamificação, a natureza dos dados pode variar amplamente, refletindo a complexidade e a diversidade dos elementos do jogo e do comportamento do usuário. Os dados a serem armazenados geralmente incluem:

\begin{itemize}
    \item \textbf{Dados de Usuários:} Informações sobre perfis de usuários, histórico de atividades, conquistas e progresso.
    \item \textbf{Dados de Atividades:} Registros de interações dos usuários com o sistema de gamificação, como pontuações, desafios concluídos e recompensas.
    \item \textbf{Dados de Feedback:} Comentários e sugestões fornecidos pelos usuários, bem como suas respostas a questionários e pesquisas.
    \item \textbf{Dados Ambientais:} Informações contextuais, como dados de sensores em ambientes físicos que podem impactar ou ser impactados pela gamificação.
\end{itemize}

Esses dados devem ser gerenciados de forma eficiente para suportar a escalabilidade e a flexibilidade necessárias em sistemas de gamificação.

\subsection{Modelos de Bancos de Dados NoSQL para Gamificação}

Os bancos de dados NoSQL são frequentemente escolhidos para aplicações de gamificação devido à sua capacidade de lidar com grandes volumes de dados não estruturados e semi-estruturados, bem como sua escalabilidade. Abaixo estão os principais modelos de bancos de dados NoSQL e suas adequações para gamificação:

\subsubsection{1. Document Stores}

\textbf{Descrição:} Bancos de dados document stores, como MongoDB e CouchDB, armazenam dados em formatos de documentos JSON ou BSON. Eles são altamente flexíveis e permitem o armazenamento de dados semi-estruturados.

\textbf{Vantagens:}
\begin{itemize}
    \item \textbf{Flexibilidade:} Adequado para armazenar dados de usuários e atividades que podem ter estruturas variadas.
    \item \textbf{Escalabilidade Horizontal:} Facilita a adição de novos nós ao cluster para lidar com o aumento de dados.
\end{itemize}

\textbf{Referência:} MongoDB oferece flexibilidade e escalabilidade que são benéficas para aplicações de gamificação onde os dados podem evoluir com o tempo \cite{chodorow2013mongodb}.

\subsubsection{2. Key-Value Stores}

\textbf{Descrição:} Bancos de dados key-value, como Redis e Riak, armazenam dados em pares de chave-valor. Eles são ideais para operações rápidas e simples, como armazenar estados de jogo e pontuações.

\textbf{Vantagens:}
\begin{itemize}
    \item \textbf{Desempenho:} Oferece alto desempenho para operações de leitura e escrita.
    \item \textbf{Simplicidade:} Ideal para armazenar e recuperar rapidamente dados simples como pontuações e rankings.
\end{itemize}

\textbf{Referência:} Redis é amplamente utilizado em cenários que exigem alta performance para manipulação de dados em tempo real, o que é essencial para a gamificação \cite{antirez2013redis}.

\subsubsection{3. Column-Family Stores}

\textbf{Descrição:} Bancos de dados column-family, como Apache Cassandra e HBase, organizam dados em colunas em vez de linhas. São adequados para dados de larga escala e consultas analíticas.

\textbf{Vantagens:}
\begin{itemize}
    \item \textbf{Escalabilidade:} Suporta grandes volumes de dados distribuídos através de múltiplos nós.
    \item \textbf{Desempenho em Consultas Analíticas:} Adequado para análises complexas e relatórios de dados de gamificação.
\end{itemize}

\textbf{Referência:} Cassandra é eficaz para aplicações que necessitam de alta disponibilidade e escalabilidade para grandes volumes de dados \cite{london2011cassandra}.

\subsection{Estratégias e Paradigmas NoSQL Mais Adequados para Gamificação}

Para escolher o banco de dados NoSQL mais adequado para uma aplicação de gamificação, considere as seguintes estratégias:

\begin{itemize}
    \item \textbf{Análise das Necessidades de Dados:} Identificar se a aplicação requer alta flexibilidade (document stores), desempenho em tempo real (key-value stores), ou escalabilidade e processamento analítico (column-family stores).
    \item \textbf{Escalabilidade e Performance:} Escolher um banco de dados que possa escalar horizontalmente para acomodar o crescimento de dados e usuários.
    \item \textbf{Modelagem de Dados:} Adaptar o modelo de dados ao tipo de banco escolhido, aproveitando suas características para otimizar a performance e a eficiência do sistema de gamificação.
\end{itemize}

\textbf{Referência:} Para uma visão mais detalhada sobre a aplicação de diferentes paradigmas NoSQL em gamificação, verifique o estudo de Xie et al. (2020) que discute como os bancos de dados NoSQL podem ser otimizados para diferentes cenários de aplicação \cite{xie2020nosql}.

\chapter{Estado da Arte}

\section{Trabalhos Relacionados a Desenvolvimento de Aplicações de Gamificação Voltadas para Eficiência Energética}

A combinação de gamificação com a eficiência energética tem sido um campo de crescente interesse, com várias abordagens sendo exploradas para incentivar comportamentos mais sustentáveis. Este capítulo revisa os principais trabalhos e estudos que abordam o desenvolvimento de aplicações de gamificação com foco na eficiência energética.

\section{Revisão de Literatura}

\subsection{Gamificação para Eficiência Energética}

A gamificação tem sido utilizada para motivar indivíduos e organizações a adotarem práticas de eficiência energética. Estudos recentes demonstram que a incorporação de elementos de jogos em sistemas de gestão de energia pode melhorar o engajamento e incentivar a redução do consumo energético. Um exemplo notável é o trabalho de \cite{del2022gamification}, que explora como a gamificação pode ser integrada em sistemas de monitoramento de energia para promover comportamentos sustentáveis.

\subsection{Estudos de Caso e Aplicações Reais}

Vários estudos de caso destacam a aplicação prática da gamificação em projetos de eficiência energética. O estudo de \cite{jones2021energy} apresenta um sistema de gamificação implementado em um ambiente corporativo, onde a competição entre equipes e recompensas foram usadas para reduzir o consumo de energia. Outro exemplo é o trabalho de \cite{smith2023smart}, que analisou a eficácia de um aplicativo de gamificação para consumidores residenciais, destacando a importância de personalizar desafios e recompensas para maximizar o impacto.

\subsection{Desafios e Limitações}

Apesar dos avanços, a implementação de gamificação para eficiência energética enfrenta vários desafios. \cite{brown2021challenges} discute as dificuldades na adaptação de elementos de jogos a contextos energéticos e os riscos de criar soluções que não sejam sustentáveis a longo prazo. A pesquisa também destaca a necessidade de equilibrar a complexidade do jogo com a facilidade de uso para garantir que as soluções sejam eficazes e adotadas amplamente.

\section{Tendências Emergentes}

\subsection{Integração com Tecnologias Emergentes}

A integração de gamificação com tecnologias emergentes, como a Internet das Coisas (IoT) e a inteligência artificial (IA), está moldando o futuro das aplicações de eficiência energética. \cite{williams2023iot} investiga como dispositivos IoT podem ser usados para fornecer feedback em tempo real e criar uma experiência de jogo mais imersiva e motivadora. Além disso, \cite{miller2024ai} explora como algoritmos de IA podem personalizar desafios e recompensas com base no comportamento do usuário, melhorando a eficácia das estratégias de gamificação.

\subsection{Personalização e Contextualização}

A personalização da experiência de gamificação para diferentes contextos e perfis de usuários é uma tendência crescente. \cite{taylor2023personalization} analisa métodos para adaptar jogos e desafios às preferências individuais e condições específicas, garantindo que a gamificação seja relevante e envolvente para cada usuário.

\section{Conclusão}

A aplicação de gamificação na eficiência energética está em constante evolução, com novas abordagens e tecnologias ampliando suas possibilidades. A revisão de trabalhos recentes destaca o potencial da gamificação para transformar comportamentos e práticas energéticas, enquanto também aponta para os desafios e oportunidades futuras. O avanço contínuo nesse campo dependerá da capacidade de integrar de forma eficaz elementos de jogos com as necessidades e contextos específicos da eficiência energética.



\chapter{Objetivo}

O objetivo geral deste trabalho é propor um modelo e uma arquitetura de dados eficiente
para a gestão de dados heterogêneos, servindo de base para uma aplicação de
gamificação. Este modelo irá integrar dados de diversas fontes, sendo identificadas:
usuários; tecnologias e sensores.
Espera-se fornecer à aplicação uma API para acesso aos serviços fornecidos pela
arquitetura proposta.
Este objetivo geral entende a realização das seguintes metas:
\begin{enumerate}
    \item Estudar e analisar as características dos bancos de dados NoSQL:
o Pesquisar e analisar os diferentes paradigmas de bancos de dados NoSQL
(baseado em documentos, chave-valor, colunar e de grafos) para identificar
o mais adequado às necessidades do projeto.
    \item Definir e implementar um modelo de dados para sensores heterogêneos:
o Desenvolver um modelo de dados flexível que possa acomodar os
diferentes formatos de dados provenientes de sensores diversos.
o Integrar princípios de gamificação na modelagem dos dados para incentivar
a redução do consumo energético.
    \item Desenvolver a interface de aplicação (API) para dispositivos móveis:
o Criar uma API baseada em nuvem que permita a comunicação entre os
dispositivos móveis e o banco de dados NoSQL.
o Garantir a segurança e a eficiência na transmissão e manipulação dos
dados pela API.
\end{enumerate}

\chapter{Metodologia}

A metodologia adotada neste projeto de Trabalho de Conclusão de Curso será
baseada em uma abordagem progressiva e interligada, que envolve uma série de
atividades contínuas e complementares. Inicialmente, será realizada uma revisão
bibliográfica abrangente para compreender as nuances e categorias dos bancos de dados
NoSQL, seguida de uma análise comparativa para determinar a melhor opção para o
contexto do projeto. Em paralelo, serão identificados os requisitos dos dados provenientes
dos sensores heterogêneos, e será projetado um modelo de dados flexível e escalável
para acomodar esses requisitos, sempre considerando os cenários de implantação que
serão levantados ao longo do projeto “Validação de sistema de gamificação, gestão e
automação na nuvem para eficiência energética de edificações públicas com foco inicial
em universidades” com a colaboração dos envolvidos. Aspectos relacionados à aplicação
de gamificação, projetada por alunos de mestrado atuantes no projeto, que utilizará esses
dados serão também levantados nesse momento e a implementação desse modelo será
acompanhada por testes de validação para garantir sua eficácia.
Para realizar a comunicação entre o que será implementado neste projeto de
conclusão e a aplicação de gamificação do projeto, será desenvolvida uma interface de
aplicação (API) para dispositivos móveis, com foco na definição de endpoints e protocolos
de comunicação que garantam a segurança e a integridade dos dados. A integração das
técnicas de gamificação será cuidadosamente incorporada à modelagem da aplicação
que gerenciará esses dados e distribuirá por meio da API, identificando pontos de
interação do usuário e implementando elementos como sistemas de recompensa e
desafios. A validação dessas técnicas será realizada por meio de testes de usabilidade e
feedback dos usuários reais.
Para validar a arquitetura e a aplicação desenvolvida, serão conduzidos testes
abrangentes de desempenho, segurança e escalabilidade. Paralelamente, será elaborada
uma documentação detalhada, incluindo relatórios técnicos que descrevem as atividades
realizadas e os resultados obtidos.

\chapter{Cronograma}

\section{Definir e implementar um modelo de dados para sensores heterogêneos}

Após a conclusão da análise das características dos bancos de dados NoSQL e a seleção do mais adequado para o projeto, o foco será na definição e implementação de um modelo de dados que possa acomodar os diferentes formatos de dados provenientes de sensores diversos. Isso incluiria:

\subsection{Identificação dos formatos de dados dos sensores}
Realizar uma análise detalhada dos diferentes tipos de dados que serão coletados pelos sensores, considerando suas variações e especificidades.

\subsection{Design do modelo de dados}
Desenvolver um modelo de dados flexível e escalável que possa acomodar os diversos formatos de dados dos sensores. Isso envolveria a definição de estruturas de dados adequadas, relacionamentos entre entidades e atributos necessários para armazenar as informações dos sensores.

\subsection{Integração dos princípios de gamificação}
Durante o processo de modelagem de dados, incorporar os princípios de gamificação para incentivar a redução do consumo energético. Isso pode envolver a definição de regras, pontos de interação e elementos de gamificação que serão integrados ao modelo de dados.

\section{Desenvolver a interface de aplicação (API) para dispositivos móveis}

Simultaneamente ao desenvolvimento do modelo de dados, será iniciado o processo de criação da interface de aplicação (API) para dispositivos móveis. Isso incluiria:

\subsection{Definição dos requisitos da API}
Identificar os requisitos funcionais e não funcionais da API, considerando aspectos como segurança, eficiência e facilidade de uso.

\subsection{Design da API}
Projetar a estrutura e os endpoints da API, determinando os métodos de comunicação e os tipos de dados que serão suportados.

\subsection{Implementação da API}
Desenvolver a API baseada em nuvem, utilizando as melhores práticas de desenvolvimento para garantir sua eficiência e segurança.

\subsection{Testes e validação da API}
Realizar testes de integração e validação da API para garantir seu funcionamento adequado e sua compatibilidade com o banco de dados NoSQL selecionado.

Após a conclusão dessas etapas, os próximos passos envolveriam a conclusão do desenvolvimento da API, a realização de testes de validação e ajustes na arquitetura como parte da Meta 4, e finalmente a documentação dos resultados do projeto e sua disseminação conforme a Meta 5.

\section{Cronograma das Atividades}

\begin{table}[h!]
\centering
\small  % or \scriptsize
\begin{tabular}{|c|c|c|c|c|c|c|c|c|c|c|c|c|}
\hline
Atividades & Jan & Fev & Mar & Abr & Mai & Jun & Jul & Ago & Set & Out & Nov & Dez \\ \hline
1.1 & X &   &   &   &   &   &   &   &   &   &   &   \\ \hline
1.2 & X & X &   &   &   &   &   &   &   &   &   &   \\ \hline
1.3 & X & X & X &   &   &   &   &   &   &   &   &   \\ \hline
1.4 & X & X & X & X &   &   &   &   &   &   &   &   \\ \hline
2.1 &   &   & X & X & X &   &   &   &   &   &   &   \\ \hline
2.2 &   &   &   & X & X & X &   &   &   &   &   &   \\ \hline
2.3 &   &   &   &   & X & X & X &   &   &   &   &   \\ \hline
3.1 &   &   &   &   &   & X & X & X &   &   &   &   \\ \hline
3.2 &   &   &   &   &   &   & X & X & X &   &   &   \\ \hline
3.3 &   &   &   &   &   &   &   & X & X & X &   &   \\ \hline
3.4 &   &   &   &   &   &   &   &   & X & X & X &   \\ \hline
4   &   &   &   &   &   &   &   &   &   & X & X & X \\ \hline
5   &   &   &   &   &   &   &   &   &   &   & X & X \\ \hline
\end{tabular}
\caption{Cronograma das Atividades}
\end{table}


\bibliographystyle{abnt}
\bibliography{bibliografia} 



\end{document}

